\documentclass[11pt]{article}
\usepackage[a4paper,margin=27mm]{geometry}
\usepackage{amsmath,amssymb,amsthm,hyperref}
\setlength{\emergencystretch}{2em}
\newtheorem{theorem}{Theorem}
\newtheorem{corollary}[theorem]{Corollary}
\newcommand{\Z}{\mathbb Z}
\title{An integral-weight obstruction to a full hyperbolic reciprocal pair}
\author{Z diophantine in Q project}
\date{19 November 2026}
\begin{document}\maketitle
\begin{abstract}
A parity certificate explains why replacing the final partial hyperbolic
pair in the 165- and 180-coordinate quartics by a full pair fails in the
retained integral Laurent compilation. The obstruction holds with any
number of stored powers of the common unit and survives every
parity-preserving change of the factor exponents. It is a restriction on
this particular compilation, not a lower bound for quartics defining the
integers.
\end{abstract}

\paragraph{Source credit.}
The shared-product constructions and the original partial hyperbolic
packing are from the notes of Gabe Conant with Astra~6.2, supplied in
\emph{H10-quartic-constructions}. The project audit and stronger extraction
are in the v78 shared-product paper~\cite{Audit}. The present parity
certificate is a subsequent project result. No new arithmetic premise is
used.

\section{The precise compilation model}
Let $c$ be an already guarded nonzero rational coordinate. An integral
Laurent compilation replaces each retained coordinate $x_i$ by
$c^{w_i}y_i$, with $w_i\in\Z$, and divides each equation by a fixed
integer power $c^{r_j}$. Powers of $c$ may be represented by additional
unit coordinates with their defining equations. The output is obtained
termwise, with each residual an ordinary polynomial of degree at most
two in the retained and unit coordinates. In particular a nonzero term
already containing two retained coordinates must have residual exponent
zero in $c$: multiplying it by an extra unit coordinate would give degree
at least three. This is the same termwise rule as in the v78 tables.

This model does not permit replacing the equations by variable-dependent
row combinations before this step, a new elliptic chart, or a nonlinear
change of the retained coordinates. Such operations lie outside the
statement below.

\begin{theorem}[Parity obstruction]\label{thm:parity}
Suppose a retained equation contains the two nonzero monomials $sU$ and
$sV$ with equal powers of $c$. Suppose another retained equation contains
$c^aUV$ and $c^bB^2$ with nonzero coefficients, where $a-b$ is odd.
These equations have no termwise quadratic integral Laurent compilation
of the specified kind, regardless of the number or powers of the added
unit coordinates.
\end{theorem}
\begin{proof}
Let $r$ be the first equation's row exponent, and let $e$ be the common
power on its two indicated terms. Quadraticity requires
\[
 e+w_s+w_U-r=0=e+w_s+w_V-r.
\]
Thus $w_U=w_V$. If $r'$ is the second row exponent, its two indicated
quadratic monomials require
\[
 a+w_U+w_V-r'=0=b+2w_B-r'.
\]
It follows that $a-b=2(w_B-w_U)$, contrary to oddness. Every indicated
monomial already has degree two, so additional stored powers of $c$
cannot alter these equalities. Fixed nonzero rational coefficients and
fixed parameter powers have weight zero and do not affect the argument.
\end{proof}

\section{The final pair of the general quartics}
The fourth shared factor has the physical terms
\[
 1+cA_1^2+c^2\overline{B_3}^{\,2}-c^3A_4^2
       +c^4\overline{B_0}^{\,4}.
\]
Here $A_1,A_4$ are raw reciprocal abscissas; they are not stored squares
of independent cleared comparison functions. The v78 presentation keeps
both and introduces one field helper $L=A_1+cA_4$. Its odd terms are
\[
 cL(A_1-cA_4)=cA_1^2-c^3A_4^2.
\]
The proposed full replacement would retain only
\[
 U=A_1+cA_4,\qquad V=A_1-cA_4,
\]
with inverse
\[
 A_1=(U+V)/2,\qquad A_4=(U-V)/(2c).
\]
This is a valid rational coordinate change when $c\ne0$. The issue is
its compatibility with the required quadratic degree.

In the retained translated point chart for the first pair,
\[
 d_1=(h_1+s_1+t_1)A_1.
\]
After the full substitution, its rational real-component row contains
$-s_{1,0}U_0/2$ and $-s_{1,0}V_0/2$ with identical unit powers.
They are distinct monomials with nonzero coefficients. In the factor's
real-component row the odd pair contributes $cU_0V_0$, whereas the
$B_3$ term contributes $c^2B_{3,0}^2$. Conjugation does not change the
coefficient of this real square. Theorem~\ref{thm:parity} applies with
$a=1$ and $b=2$.

\begin{corollary}
The full replacement of this raw pair cannot eliminate its two partial
helper coordinates by integral diagonal Laurent weights in the retained
quadratic equations. The conclusion applies to both general quartics,
and to both the original and strengthened v78 presentations. It also
applies when the odd and even factor exponents are changed by arbitrary
even integers while the indicated point row is retained.
\end{corollary}
\begin{proof}
The two constructions share these point and fourth-factor equations.
The v78 strengthening changes only the first three factors. More general
parity-preserving exponent changes retain the hypotheses of the theorem.
\end{proof}

The same proof applies to a proposed full pair involving any raw
reciprocal whose retained point equation supplies the two indicated
bilinear monomials. This condition must be checked in each application;
it is not asserted for arbitrary elliptic charts or every possible
choice of detector.

\section{An exact certificate and its limitations}
The independently runnable checker
\begin{center}\small
\nolinkurl{computation/variable_reductions_v79/quartics/}\par
\nolinkurl{check_hyperbolic_parity.py}
\end{center}
reads the pinned Laurent source and actual fourth-factor assignment,
checks the two source monomials and their substitution, and verifies
$c((U+V)/2)^2-c^3((U-V)/(2c))^2=cUV$ coefficient by coefficient.
It records the four integer-weight equations in unknowns
$(w_s,w_U,w_V,w_B,r,r')$. Adding all four equations modulo two gives
$0=1$. This is a complete finite certificate; it uses neither a numerical
search range nor an optimizer's infeasibility report.

The result explains a failure of a particular algebraic compression.
It proves no global minimum of 165 or 180 variables, and does not
exclude smaller quartics from other charts, other gate equations,
nonlinear transformations, or different arithmetic constructions.
No variable-count improvement is inferred from this obstruction.

\begin{thebibliography}{9}
\bibitem{Audit} Z diophantine in Q project, \emph{Audit and extension of the
shared-product quartics}, 19 November 2026,
\nolinkurl{papers/quartic/shared_product_audit_20261119_v78.tex}.
\bibitem{Chart} Z diophantine in Q project, \emph{A 193-coordinate quartic},
v74, \nolinkurl{papers/quartic/quartic193_20261117_v74.tex}.
\end{thebibliography}
\end{document}
